\part{Distribution et installation}
\label{part:install}

%% On compte sur la commande path du package url
\newcommand{\pathfmt}[1]{\path{#1}}
\newcommand{\clavier}[1]{\texttt{#1}}

  Les bibliothèques de fonctions qui sont au c\oe{}ur du système AIDA
(manipulation de fichiers, traitements) sont écrites en langage C et en
langage Python.

\section{Pré-requis}
  Pour installer et utiliser les différents programmes qui constituent
AIDA sur une machine, celle-ci doit offrir un certain nombre de
services. Ce sont :

\begin{description}
  \item[un système de développement \texttt{C-ANSI}] : un compilateur, un
    éditeur de liens, les programmes \texttt{make} et \texttt{makedepend}.
    AIDA est développé avec les versions standard disponibles sur la
    distribution GNU/Linux \textbf{Ubuntu} (architecture x86\_64),
    avec gcc et clang.

  \item[python] : les programmes python dans AIDA sont compatibles avec la
    version installée en standard sur Ubuntu. La présence des packages
    nécessaires est testée par l'installateur de AIDA.

  \item[\LaTeX] : cette documentation, et les documentations des fonctions de
    traitement («opérateurs») sont écrites en \LaTeX.

  \item[la bibliothèque udunits2] : c'est une bibliothèque de manipulation
    d'unités développée par le groupe \textit{UNIDATA}. Dans AIDA,
    elle sert, par exemple, à vérifier que l'unité d'un vecteur correspond
    bien à ce que le programme de traitement attend.

  \item[la bibliothèque lapack (et blas)] : Lapack (\textbf{L}inear
    \textbf{A}lgebra \textbf{PACK}age), est une bibliothèque de routines
    Fortran 77 utilisée pour résoudre des problèmes d'algèbre
    linéaire : résolution de systèmes linéaires, problèmes de moindres carrés,
    problèmes aux valeurs propres, factorisations de matrice, ...

  \item[la bibliothèque Z] : la bibliothèque zlib permet de compresser ou de
    décompresser des données en mémoire ou de lire et écrire des fichiers au
    format gzip.

  \item[cextract] : ce programme génère automatiquement les prototypes de
    fonctions écrites en langage C. Ce programme n'est donc nécessaire
    \emph{que lors de modifications des sources}. Il n'est pas nécessaires
    pour une utilisation courante. L'auteur est Adam Bryant (adb@bu.edu). AIDA
    utilise la version 1.7, qui date de 1992.
    Ce programme qui a été diffusé sur le newsgroup
    \texttt{comp}\texttt{.sources.reviewed}, et que l'on pouvait trouver à
    \url{http://www.cs.bu.edu/ftp/pub/adb/}, ou inclus dans divers packages
    logiciels, par exemple les sources de Thot
    (\url{http://www.inrialpes.fr/opera/Thot/index.html}) a été enregistré
    dans \$AIDA\_HOME/extern/src/.  Il est accompagné de
    \texttt{cextdoc}, générateur automatique de documentation, qui n'est plus
    utilisé dans AIDA.
    Si les difficultés d'utilisation deviennent trop grandes,
    une alternative serait Makeheaders (\url{http://www.hwaci.com/sw/mkhdr/})

  \item[doxygen] : doxygen est un logiciel permettant de générer de la
    documentation technique à partir du code source d'un programme.
    Il extrait les déclarations et certains commentaire du code source,
    et crée un manuel de référence (format HTML, ou pdf via LaTeX),
    pour permettre au programmeur de localiser rapidement une fonction
    ou une structure quelconque. C'est un logiciel libre.

\end{description}

D'autres outils peuvent être nécessaires pour des applications qui ne sont pas
utilisées partout :

\begin{description}
  \item[Labview] : LabVIEW est un logiciel de développement d'applications
    par programmation graphique. Il est développé et vendu par la
    société américaine National Instruments (\url{http://www.ni.com/labview/f/}).
    Il est principalement orienté vers les tests, la mesure et
    l'automatisation, mais il permet aussi de développer des applications
    d'analyse de données.

  \item[la bibliothèque netcdf] : elle sert à lire et écrire des fichiers au
    format NetCDF, et ne sera utile que si vous avez à manipuler des fichiers
    à ce format.

  \item[graphviz] : logiciel de visualisation de graphes
    \url{http://www.graphviz.org/}.

\end{description}


\section{Modes de distributions et installation}

La méthode de distribution de AIDA est basée sur un fichier \texttt{tar}.
%%%Se repporter à la documentation \texttt{\textbf{installateur\_aida.pdf}}.
Se repporter à la documentation \url{installateur\_aida.pdf}.

